\documentclass[10pt,a4paper]{amsart}
\usepackage{amsmath,amssymb,amsthm,graphicx}
% Retained minimal preamble for the Stacks Project chapter more-algebra.tex (More on Algebra).
% The chapter file pinned-original-source.tex carries only \input{preamble} on its first line, so no
% preamble text lives inside the chapter itself.  The Stacks Project uses one single shared preamble for
% every chapter; the snapshot stages it identically as sources/stacks-chow/preamble.tex and
% sources/stacks-etale-cohomology/preamble.tex.  The block below is copied mechanically, line for line,
% from sources/stacks-chow/preamble.tex lines 14-18 (the Xy-pic setup) and 154-217 (the native theorem
% declarations together with the mathematical macros).  Nothing was added, removed or redefined; the
% wrapper class and the amsmath/amssymb/amsthm/graphicx base packages are supplied by the assembler.
% Chapter: more-algebra.tex -- More on Algebra.
% For commutative diagrams we use Xy-pic
\usepackage[all]{xy}

% We use 2cell for 2-commutative diagrams.
\xyoption{2cell}

% Theorem environments.
%
\theoremstyle{plain}
\newtheorem{theorem}[subsection]{Theorem}
\newtheorem{proposition}[subsection]{Proposition}
\newtheorem{lemma}[subsection]{Lemma}

\theoremstyle{definition}
\newtheorem{definition}[subsection]{Definition}
\newtheorem{example}[subsection]{Example}
\newtheorem{exercise}[subsection]{Exercise}
\newtheorem{situation}[subsection]{Situation}

\theoremstyle{remark}
\newtheorem{remark}[subsection]{Remark}
\newtheorem{remarks}[subsection]{Remarks}

\numberwithin{equation}{subsection}

% Macros
%
\def\lim{\mathop{\mathrm{lim}}\nolimits}
\def\colim{\mathop{\mathrm{colim}}\nolimits}
\def\Spec{\mathop{\mathrm{Spec}}}
\def\Hom{\mathop{\mathrm{Hom}}\nolimits}
\def\Ext{\mathop{\mathrm{Ext}}\nolimits}
\def\SheafHom{\mathop{\mathcal{H}\!\mathit{om}}\nolimits}
\def\SheafExt{\mathop{\mathcal{E}\!\mathit{xt}}\nolimits}
\def\Sch{\mathit{Sch}}
\def\Mor{\mathop{\mathrm{Mor}}\nolimits}
\def\Ob{\mathop{\mathrm{Ob}}\nolimits}
\def\Sh{\mathop{\mathit{Sh}}\nolimits}
\def\NL{\mathop{N\!L}\nolimits}
\def\CH{\mathop{\mathrm{CH}}\nolimits}
\def\proetale{{pro\text{-}\acute{e}tale}}
\def\etale{{\acute{e}tale}}
\def\QCoh{\mathit{QCoh}}
\def\Ker{\mathop{\mathrm{Ker}}}
\def\Im{\mathop{\mathrm{Im}}}
\def\Coker{\mathop{\mathrm{Coker}}}
\def\Coim{\mathop{\mathrm{Coim}}}

% Boxtimes
%
\DeclareMathSymbol{\boxtimes}{\mathbin}{AMSa}{"02}

%
% Macros for moduli stacks/spaces
%
\def\QCohstack{\mathcal{QC}\!\mathit{oh}}
\def\Cohstack{\mathcal{C}\!\mathit{oh}}
\def\Spacesstack{\mathcal{S}\!\mathit{paces}}
\def\Quotfunctor{\mathrm{Quot}}
\def\Hilbfunctor{\mathrm{Hilb}}
\def\Curvesstack{\mathcal{C}\!\mathit{urves}}
\def\Polarizedstack{\mathcal{P}\!\mathit{olarized}}
\def\Complexesstack{\mathcal{C}\!\mathit{omplexes}}
% \Pic is the operator that assigns to X its picard group, usage \Pic(X)
% \Picardstack_{X/B} denotes the Picard stack of X over B
% \Picardfunctor_{X/B} denotes the Picard functor of X over B
\def\Pic{\mathop{\mathrm{Pic}}\nolimits}
\def\Picardstack{\mathcal{P}\!\mathit{ic}}
\def\Picardfunctor{\mathrm{Pic}}
\def\Deformationcategory{\mathcal{D}\!\mathit{ef}}

% --- appended 2026-10-12 (Chief): the project own definitions of the reference/slogan/history
% display environments, copied mechanically from the shared Stacks preamble lines 8-12
% (\usepackage{verbatim} plus the three \newenvironment lines, which the chapter preamble keeps
% verbatim from the project preamble). Nothing added, removed or redefined: these environments
% carry no mathematical content and the project itself typesets them as comments. Without these
% three lines a counted statement whose own text uses them fails with Environment slogan undefined.
\usepackage{verbatim}
\newenvironment{reference}{\comment}{\endcomment}
\newenvironment{slogan}{\comment}{\endcomment}
\newenvironment{history}{\comment}{\endcomment}
\usepackage{cleveref}
\crefname{originallemma}{original lemma}{original lemmas}
\Crefname{originallemma}{Original lemma}{Original lemmas}
\makeatletter
\newlabel{algebra-lemma-catenary-check-irreducible}{{0AUP}{source}{}{sourcepointer-1}{}}
\newlabel{algebra-lemma-catenary-check-irreducible@cref}{{[originallemma][0][]0AUP}{[1][0][]source}{}{}{}}
\makeatother
\crefname{originallemma}{original lemma}{original lemmas}
\Crefname{originallemma}{Original lemma}{Original lemmas}
\makeatletter
\newlabel{algebra-lemma-catenary-check-local}{{0AUN}{source}{}{sourcepointer-2}{}}
\newlabel{algebra-lemma-catenary-check-local@cref}{{[originallemma][0][]0AUN}{[1][0][]source}{}{}{}}
\makeatother
\crefname{originallemma}{original lemma}{original lemmas}
\Crefname{originallemma}{Original lemma}{Original lemmas}
\makeatletter
\newlabel{algebra-lemma-completion-flat}{{00MB}{source}{}{sourcepointer-3}{}}
\newlabel{algebra-lemma-completion-flat@cref}{{[originallemma][0][]00MB}{[1][0][]source}{}{}{}}
\makeatother
\crefname{originallemma}{original lemma}{original lemmas}
\Crefname{originallemma}{Original lemma}{Original lemmas}
\makeatletter
\newlabel{algebra-lemma-dimension-base-fibre-equals-total}{{00ON}{source}{}{sourcepointer-4}{}}
\newlabel{algebra-lemma-dimension-base-fibre-equals-total@cref}{{[originallemma][0][]00ON}{[1][0][]source}{}{}{}}
\makeatother
\crefname{originallemma}{original lemma}{original lemmas}
\Crefname{originallemma}{Original lemma}{Original lemmas}
\makeatletter
\newlabel{algebra-lemma-localization-at-closed-point-special-fibre}{{0AUG}{source}{}{sourcepointer-5}{}}
\newlabel{algebra-lemma-localization-at-closed-point-special-fibre@cref}{{[originallemma][0][]0AUG}{[1][0][]source}{}{}{}}
\makeatother
\crefname{originallemma}{original lemma}{original lemmas}
\Crefname{originallemma}{Original lemma}{Original lemmas}
\makeatletter
\newlabel{algebra-lemma-localization-catenary}{{00NJ}{source}{}{sourcepointer-6}{}}
\newlabel{algebra-lemma-localization-catenary@cref}{{[originallemma][0][]00NJ}{[1][0][]source}{}{}{}}
\makeatother
\crefname{originallemma}{original lemma}{original lemmas}
\Crefname{originallemma}{Original lemma}{Original lemmas}
\makeatletter
\newlabel{algebra-lemma-quotient-catenary}{{00NK}{source}{}{sourcepointer-7}{}}
\newlabel{algebra-lemma-quotient-catenary@cref}{{[originallemma][0][]00NK}{[1][0][]source}{}{}{}}
\makeatother
\crefname{originallemma}{original lemma}{original lemmas}
\Crefname{originallemma}{Original lemma}{Original lemmas}
\makeatletter
\newlabel{algebra-remark-Noetherian-complete-local-ring-universally-catenary}{{032C}{source}{}{sourcepointer-8}{}}
\newlabel{algebra-remark-Noetherian-complete-local-ring-universally-catenary@cref}{{[originallemma][0][]032C}{[1][0][]source}{}{}{}}
\makeatother
\crefname{originallemma}{original lemma}{original lemmas}
\Crefname{originallemma}{Original lemma}{Original lemmas}
\makeatletter
\newlabel{lemma-flat-under-catenary-equidimensional}{{32389}{source}{}{sourcepointer-9}{}}
\newlabel{lemma-flat-under-catenary-equidimensional@cref}{{[originallemma][0][]32389}{[1][0][]source}{}{}{}}
\makeatother
\title{A formally catenary Noetherian local ring is universally catenary}
\author{The Stacks Project Authors}
\date{stacks-project:more-algebra.tex:local-snapshot-20261012; GNU FDL 1.2 or later}
\begin{document}
\maketitle
\setcounter{subsection}{1}
\noindent\textit{Editorial note.} All delivered mathematical text is the author's own, copied from the pinned Stacks Project chapter file more-algebra.tex; no statement or proof step was rewritten, reordered or removed.\par\medskip
\noindent\textit{Editorial note.} The chapter file pinned-original-source.tex carries only \textbackslash{}input\{preamble\} on its first line, so no preamble text lives inside the chapter. The retained file sources/stacks-more-algebra/retained-preamble.tex is a mechanical, line-for-line copy of the project's single shared preamble as staged in the snapshot; nothing was added, removed or redefined and no package was added.\par\medskip
\noindent\textit{Editorial note.} The source URL is the source-level Stacks Project URL, as required when more than one card is produced from one source; the per-claim tag is carried by claim\_key.\par\medskip
\noindent\textit{Editorial note.} The countable proof span is the author's proof environment at source lines 32443-32481. The nine references inside it are recorded as non-load-bearing navigation pointers; they name standard facts about catenarity, quotients, localizations and completion whose applications stay inside the counted proof span.\par\medskip
\noindent\textit{Allowed standard prerequisites (editorial).}\begin{itemize}
\item Catenary and universally catenary rings, and the criterion that every local ring essentially of finite type be catenary. used at proof lines 32444-32449. Is the notion being established and the reduction that turns it into a local statement about polynomial extensions.
\item Formal catenarity and equidimensionality of the spectrum of a completion. the hypothesis at statement line 32439 and used at proof lines 32446 and 32462. Is the input that settles the completed polynomial ring.
\item Flatness of the completion map and the existence of a unique maximal ideal of the completed polynomial ring over a given one. used at proof lines 32456-32461. Is what lets the catenarity question be transferred to the completion.
\end{itemize}\medskip
\ifdefined\hypertarget\hypertarget{sourcepointer-1}{}\fi
\noindent\textit{Source reference.} Original lemma 0AUP: focused-seeds/AGP/sources/stacks-snapshot-20261011/tags.txt, source lines 13657--13657. Cross-chapter navigation pointer into the Stacks Project; the standard fact it supplies is applied inside the counted proof span and the argument resting on it stays there.\par\medskip
\ifdefined\hypertarget\hypertarget{sourcepointer-2}{}\fi
\noindent\textit{Source reference.} Original lemma 0AUN: focused-seeds/AGP/sources/stacks-snapshot-20261011/tags.txt, source lines 13656--13656. Cross-chapter navigation pointer into the Stacks Project; the standard fact it supplies is applied inside the counted proof span and the argument resting on it stays there.\par\medskip
\ifdefined\hypertarget\hypertarget{sourcepointer-3}{}\fi
\noindent\textit{Source reference.} Original lemma 00MB: focused-seeds/AGP/sources/stacks-snapshot-20261011/tags.txt, source lines 817--817. Cross-chapter navigation pointer into the Stacks Project; the standard fact it supplies is applied inside the counted proof span and the argument resting on it stays there.\par\medskip
\ifdefined\hypertarget\hypertarget{sourcepointer-4}{}\fi
\noindent\textit{Source reference.} Original lemma 00ON: focused-seeds/AGP/sources/stacks-snapshot-20261011/tags.txt, source lines 901--901. Cross-chapter navigation pointer into the Stacks Project; the standard fact it supplies is applied inside the counted proof span and the argument resting on it stays there.\par\medskip
\ifdefined\hypertarget\hypertarget{sourcepointer-5}{}\fi
\noindent\textit{Source reference.} Original lemma 0AUG: focused-seeds/AGP/sources/stacks-snapshot-20261011/tags.txt, source lines 13649--13649. Cross-chapter navigation pointer into the Stacks Project; the standard fact it supplies is applied inside the counted proof span and the argument resting on it stays there.\par\medskip
\ifdefined\hypertarget\hypertarget{sourcepointer-6}{}\fi
\noindent\textit{Source reference.} Original lemma 00NJ: focused-seeds/AGP/sources/stacks-snapshot-20261011/tags.txt, source lines 861--861. Cross-chapter navigation pointer into the Stacks Project; the standard fact it supplies is applied inside the counted proof span and the argument resting on it stays there.\par\medskip
\ifdefined\hypertarget\hypertarget{sourcepointer-7}{}\fi
\noindent\textit{Source reference.} Original lemma 00NK: focused-seeds/AGP/sources/stacks-snapshot-20261011/tags.txt, source lines 862--862. Cross-chapter navigation pointer into the Stacks Project; the standard fact it supplies is applied inside the counted proof span and the argument resting on it stays there.\par\medskip
\ifdefined\hypertarget\hypertarget{sourcepointer-8}{}\fi
\noindent\textit{Source reference.} Original lemma 032C: focused-seeds/AGP/sources/stacks-snapshot-20261011/tags.txt, source lines 4015--4015. Cross-chapter navigation pointer into the Stacks Project; the standard fact it supplies is applied inside the counted proof span and the argument resting on it stays there.\par\medskip
\ifdefined\hypertarget\hypertarget{sourcepointer-9}{}\fi
\noindent\textit{Source reference.} Original lemma 32389: sources/stacks-more-algebra/pinned-original-source.tex, source lines 32389--32389. Same-chapter navigation pointer; the statement it names is standard for the section and the argument that consumes it stays inside the counted proof span.\par\medskip
\begin{lemma}
\label{lemma-formally-catenary}
Let $A$ be a formally catenary Noetherian local ring.
Then $A$ is universally catenary.
\end{lemma}

\begin{proof}
We may replace $A$ by $A/\mathfrak p$ where $\mathfrak p$ is a minimal prime
of $A$, see Algebra, Lemma \ref{algebra-lemma-catenary-check-irreducible}.
Thus we may assume that the spectrum of $A^\wedge$ is equidimensional.
It suffices to show that every local ring essentially of finite type
over $A$ is catenary (see for example
Algebra, Lemma \ref{algebra-lemma-catenary-check-local}).
Hence it suffices to show that $A[x_1, \ldots, x_n]_\mathfrak m$ is catenary
where $\mathfrak m \subset A[x_1, \ldots, x_n]$ is a maximal
ideal lying over $\mathfrak m_A$, see
Algebra, Lemma \ref{algebra-lemma-localization-at-closed-point-special-fibre}
(and Algebra, Lemmas \ref{algebra-lemma-quotient-catenary} and
\ref{algebra-lemma-localization-catenary}).
Let $\mathfrak m' \subset A^\wedge[x_1, \ldots, x_n]$ be the unique
maximal ideal lying over $\mathfrak m$. Then
$$
A[x_1, \ldots, x_n]_\mathfrak m \to A^\wedge[x_1, \ldots, x_n]_{\mathfrak m'}
$$
is local and flat (Algebra, Lemma \ref{algebra-lemma-completion-flat}).
Hence it suffices to show that the ring on the right
hand side catenary with equidimensional spectrum, see
Lemma \ref{lemma-flat-under-catenary-equidimensional}.
It is catenary because complete local rings are universally catenary
(Algebra, Remark
\ref{algebra-remark-Noetherian-complete-local-ring-universally-catenary}).
Pick any minimal prime $\mathfrak q$ of
$A^\wedge[x_1, \ldots, x_n]_{\mathfrak m'}$. Then
$\mathfrak q = \mathfrak p A^\wedge[x_1, \ldots, x_n]_{\mathfrak m'}$
for some minimal prime $\mathfrak p$ of $A^\wedge$ (small detail omitted).
Hence
$$
\dim(A^\wedge[x_1, \ldots, x_n]_{\mathfrak m'}/\mathfrak q) =
\dim(A^\wedge/\mathfrak p) + n = \dim(A^\wedge) + n
$$
the first equality by
Algebra, Lemma \ref{algebra-lemma-dimension-base-fibre-equals-total}
and the second because the spectrum of $A^\wedge$ is equidimensional.
This finishes the proof.
\end{proof}

\end{document}
